\documentclass[10pt,a4paper]{amsart}
\usepackage{amsmath,amssymb,amsthm,graphicx}
\usepackage{verbatim}
\newenvironment{reference}{\comment}{\endcomment}
%\newenvironment{reference}{}{}
\newenvironment{slogan}{\comment}{\endcomment}
\newenvironment{history}{\comment}{\endcomment}

% For commutative diagrams we use Xy-pic
\usepackage[all]{xy}

% We use 2cell for 2-commutative diagrams.
\xyoption{2cell}
% Theorem environments.
%
\theoremstyle{plain}
\newtheorem{theorem}[subsection]{Theorem}
\newtheorem{proposition}[subsection]{Proposition}
\newtheorem{lemma}[subsection]{Lemma}

\theoremstyle{definition}
\newtheorem{definition}[subsection]{Definition}
\newtheorem{example}[subsection]{Example}
\newtheorem{exercise}[subsection]{Exercise}
\newtheorem{situation}[subsection]{Situation}

\theoremstyle{remark}
\newtheorem{remark}[subsection]{Remark}
\newtheorem{remarks}[subsection]{Remarks}

\numberwithin{equation}{subsection}

% Macros
%
\def\lim{\mathop{\mathrm{lim}}\nolimits}
\def\colim{\mathop{\mathrm{colim}}\nolimits}
\def\Spec{\mathop{\mathrm{Spec}}}
\def\Hom{\mathop{\mathrm{Hom}}\nolimits}
\def\Ext{\mathop{\mathrm{Ext}}\nolimits}
\def\SheafHom{\mathop{\mathcal{H}\!\mathit{om}}\nolimits}
\def\SheafExt{\mathop{\mathcal{E}\!\mathit{xt}}\nolimits}
\def\Sch{\mathit{Sch}}
\def\Mor{\mathop{\mathrm{Mor}}\nolimits}
\def\Ob{\mathop{\mathrm{Ob}}\nolimits}
\def\Sh{\mathop{\mathit{Sh}}\nolimits}
\def\NL{\mathop{N\!L}\nolimits}
\def\CH{\mathop{\mathrm{CH}}\nolimits}
\def\proetale{{pro\text{-}\acute{e}tale}}
\def\etale{{\acute{e}tale}}
\def\QCoh{\mathit{QCoh}}
\def\Ker{\mathop{\mathrm{Ker}}}
\def\Im{\mathop{\mathrm{Im}}}
\def\Coker{\mathop{\mathrm{Coker}}}
\def\Coim{\mathop{\mathrm{Coim}}}

% Boxtimes
%
\DeclareMathSymbol{\boxtimes}{\mathbin}{AMSa}{"02}

%
% Macros for moduli stacks/spaces
%
\def\QCohstack{\mathcal{QC}\!\mathit{oh}}
\def\Cohstack{\mathcal{C}\!\mathit{oh}}
\def\Spacesstack{\mathcal{S}\!\mathit{paces}}
\def\Quotfunctor{\mathrm{Quot}}
\def\Hilbfunctor{\mathrm{Hilb}}
\def\Curvesstack{\mathcal{C}\!\mathit{urves}}
\def\Polarizedstack{\mathcal{P}\!\mathit{olarized}}
\def\Complexesstack{\mathcal{C}\!\mathit{omplexes}}
% \Pic is the operator that assigns to X its picard group, usage \Pic(X)
% \Picardstack_{X/B} denotes the Picard stack of X over B
% \Picardfunctor_{X/B} denotes the Picard functor of X over B
\def\Pic{\mathop{\mathrm{Pic}}\nolimits}
\def\Picardstack{\mathcal{P}\!\mathit{ic}}
\def\Picardfunctor{\mathrm{Pic}}
\def\Deformationcategory{\mathcal{D}\!\mathit{ef}}
\usepackage{cleveref}
\crefname{originallemma}{original lemma}{original lemmas}
\Crefname{originallemma}{Original lemma}{Original lemmas}
\makeatletter
\newlabel{algebra-lemma-bound-depth}{{00LK}{source}{}{sourcepointer-1}{}}
\newlabel{algebra-lemma-bound-depth@cref}{{[originallemma][0][]00LK}{[1][0][]source}{}{}{}}
\makeatother
\crefname{originallemma}{original lemma}{original lemmas}
\Crefname{originallemma}{Original lemma}{Original lemmas}
\makeatletter
\newlabel{algebra-lemma-completion-generalities}{{0315}{source}{}{sourcepointer-2}{}}
\newlabel{algebra-lemma-completion-generalities@cref}{{[originallemma][0][]0315}{[1][0][]source}{}{}{}}
\makeatother
\crefname{originallemma}{original lemma}{original lemmas}
\Crefname{originallemma}{Original lemma}{Original lemmas}
\makeatletter
\newlabel{algebra-lemma-criterion-reduced}{{031R}{source}{}{sourcepointer-3}{}}
\newlabel{algebra-lemma-criterion-reduced@cref}{{[originallemma][0][]031R}{[1][0][]source}{}{}{}}
\makeatother
\crefname{originallemma}{original lemma}{original lemmas}
\Crefname{originallemma}{Original lemma}{Original lemmas}
\makeatletter
\newlabel{algebra-lemma-regular-graded}{{00NO}{source}{}{sourcepointer-4}{}}
\newlabel{algebra-lemma-regular-graded@cref}{{[originallemma][0][]00NO}{[1][0][]source}{}{}{}}
\makeatother
\crefname{originalreference}{original reference}{original references}
\Crefname{originalreference}{Original reference}{Original references}
\makeatletter
\newlabel{lemma-reduced-quotient-regular-ring-dim-2}{{0C48}{source}{}{sourcepointer-5}{}}
\newlabel{lemma-reduced-quotient-regular-ring-dim-2@cref}{{[originalreference][0][]0C48}{[1][0][]source}{}{}{}}
\makeatother
\crefname{originallemma}{original lemma}{original lemmas}
\Crefname{originallemma}{Original lemma}{Original lemmas}
\makeatletter
\newlabel{more-algebra-lemma-completion-depth}{{07NW}{source}{}{sourcepointer-6}{}}
\newlabel{more-algebra-lemma-completion-depth@cref}{{[originallemma][0][]07NW}{[1][0][]source}{}{}{}}
\makeatother
\crefname{originallemma}{original lemma}{original lemmas}
\Crefname{originallemma}{Original lemma}{Original lemmas}
\makeatletter
\newlabel{more-algebra-lemma-lift-residue-field}{{0C34}{source}{}{sourcepointer-7}{}}
\newlabel{more-algebra-lemma-lift-residue-field@cref}{{[originallemma][0][]0C34}{[1][0][]source}{}{}{}}
\makeatother
\title{Noetherian local k-algebras whose completed tensor is a nondegenerate quadric power series ring: three equivalent characterizations}
\author{The Stacks Project Authors}
\date{stacks-project:curves.tex:local-snapshot-20261012; GNU FDL 1.2 or later}
\begin{document}
\maketitle
\noindent\textit{Editorial note.} All delivered mathematical text is the author's own, copied from the pinned Stacks Project chapter file curves.tex; no statement or proof step was rewritten, reordered or removed.\par\medskip
\noindent\textit{Editorial note.} The counted proof is the author's complete argument: the reduction of (3) to (1) via irreducibility of the quadric, the exclusion argument for dimension, the full lifting induction with the congruence modulo m\textasciicircum{}4 and the nondegeneracy basis step, the completeness limit, and the kernel identification with its degree contradiction are all retained verbatim.\par\medskip
\noindent\textit{Editorial note.} The label lemma-nodal-algebraic names the counted statement, and the proof environment immediately following it in the source is delivered verbatim as its proof.\par\medskip
\noindent\textit{Editorial note.} The counted statement is carried by the author's own lemma environment in the chapter source.\par\medskip
\noindent\textit{Allowed standard prerequisites (editorial).}\begin{itemize}
\item Depth and dimension estimates on local rings. Algebra, Lemmas algebra-lemma-criterion-reduced, algebra-lemma-regular-graded, algebra-lemma-bound-depth. Rules out the Artinian and the dimension 2 (regular) cases in the implication (1) => (2), forcing depth(A) = 1.
\item Lifting residue fields and completion generalities. More on Algebra, Lemmas more-algebra-lemma-completion-depth and more-algebra-lemma-lift-residue-field; Algebra, Lemma algebra-lemma-completion-generalities. Lets one assume A complete with a coefficient field kappa, so that the successive improvement of generators converges.
\end{itemize}\medskip
\ifdefined\hypertarget\hypertarget{sourcepointer-1}{}\fi
\noindent\textit{Source reference.} Original lemma 00LK: sources/stacks-chow/tags.txt, source lines 790--790. Prerequisite pointer invoked at proof line(s) 4123; it names a standard fact of a sibling chapter that the counted argument uses, and every step of the counted proof built on it is retained verbatim in the delivered body.\par\medskip
\ifdefined\hypertarget\hypertarget{sourcepointer-2}{}\fi
\noindent\textit{Source reference.} Original lemma 0315: sources/stacks-chow/tags.txt, source lines 3972--3972. Prerequisite pointer invoked at proof line(s) 4171; it names a standard fact of a sibling chapter that the counted argument uses, and every step of the counted proof built on it is retained verbatim in the delivered body.\par\medskip
\ifdefined\hypertarget\hypertarget{sourcepointer-3}{}\fi
\noindent\textit{Source reference.} Original lemma 031R: sources/stacks-chow/tags.txt, source lines 3994--3994. Prerequisite pointer invoked at proof line(s) 4118; it names a standard fact of a sibling chapter that the counted argument uses, and every step of the counted proof built on it is retained verbatim in the delivered body.\par\medskip
\ifdefined\hypertarget\hypertarget{sourcepointer-4}{}\fi
\noindent\textit{Source reference.} Original lemma 00NO: sources/stacks-chow/tags.txt, source lines 866--866. Prerequisite pointer invoked at proof line(s) 4121; it names a standard fact of a sibling chapter that the counted argument uses, and every step of the counted proof built on it is retained verbatim in the delivered body.\par\medskip
\ifdefined\hypertarget\hypertarget{sourcepointer-5}{}\fi
\noindent\textit{Source reference.} Original reference 0C48: sources/stacks-curves/pinned-original-source.tex, source lines 4034--4034. Prerequisite pointer invoked at proof line(s) 4172; it names a standard fact of the same chapter that the counted argument uses, and every step of the counted proof built on it is retained verbatim in the delivered body.\par\medskip
\ifdefined\hypertarget\hypertarget{sourcepointer-6}{}\fi
\noindent\textit{Source reference.} Original lemma 07NW: sources/stacks-chow/tags.txt, source lines 9776--9776. Prerequisite pointer invoked at proof line(s) 4128; it names a standard fact of a sibling chapter that the counted argument uses, and every step of the counted proof built on it is retained verbatim in the delivered body.\par\medskip
\ifdefined\hypertarget\hypertarget{sourcepointer-7}{}\fi
\noindent\textit{Source reference.} Original lemma 0C34: sources/stacks-chow/tags.txt, source lines 15179--15179. Prerequisite pointer invoked at proof line(s) 4131; it names a standard fact of a sibling chapter that the counted argument uses, and every step of the counted proof built on it is retained verbatim in the delivered body.\par\medskip
\begin{lemma}
\label{lemma-nodal-algebraic}
Let $k$ be a field. Let $(A, \mathfrak m, \kappa)$ be a Noetherian
local $k$-algebra. The following are equivalent
\begin{enumerate}
\item $\kappa/k$ is separable, $A$ is reduced,
$\dim_\kappa(\mathfrak m/\mathfrak m^2) = 2$, and there exists a nondegenerate
$q \in \text{Sym}^2_\kappa(\mathfrak m/\mathfrak m^2)$
which maps to zero in $\mathfrak m^2/\mathfrak m^3$,
\item $\kappa/k$ is separable, $\text{depth}(A) = 1$,
$\dim_\kappa(\mathfrak m/\mathfrak m^2) = 2$, and there exists a nondegenerate
$q \in \text{Sym}^2_\kappa(\mathfrak m/\mathfrak m^2)$
which maps to zero in $\mathfrak m^2/\mathfrak m^3$,
\item $\kappa/k$ is separable,
$A^\wedge \cong \kappa[[x, y]]/(ax^2 + bxy + cy^2)$
as a $k$-algebra where $ax^2 + bxy + cy^2$ is a nondegenerate quadratic form
over $\kappa$.
\end{enumerate}
\end{lemma}

\begin{proof}
Assume (3). Then $A^\wedge$ is reduced because $ax^2 + bxy + cy^2$
is either irreducible or a product of two nonassociated prime elements.
Hence $A \subset A^\wedge$ is reduced. It follows that (1) is true.

\medskip\noindent
Assume (1). Then $A$ cannot be Artinian, since it would not be reduced
because $\mathfrak m \not = (0)$.
Hence $\dim(A) \geq 1$, hence $\text{depth}(A) \geq 1$
by Algebra, Lemma \ref{algebra-lemma-criterion-reduced}.
On the other hand $\dim(A) = 2$ implies $A$ is regular
which contradicts the existence of $q$ by
Algebra, Lemma \ref{algebra-lemma-regular-graded}.
Thus $\dim(A) \leq 1$ and we conclude $\text{depth}(A) = 1$
by Algebra, Lemma \ref{algebra-lemma-bound-depth}.
It follows that (2) is true.

\medskip\noindent
Assume (2). Since the depth of $A$ is the same as the depth of $A^\wedge$
(More on Algebra, Lemma \ref{more-algebra-lemma-completion-depth})
and since the other conditions are insensitive to completion, we may
assume that $A$ is complete. Choose $\kappa \to A$ as in
More on Algebra, Lemma \ref{more-algebra-lemma-lift-residue-field}.
Since $\dim_\kappa(\mathfrak m/\mathfrak m^2) = 2$ we can choose
$x_0, y_0 \in \mathfrak m$ which map to a basis.
We obtain a continuous $\kappa$-algebra map
$$
\kappa[[x, y]] \longrightarrow A
$$
by the rules $x \mapsto x_0$ and $y \mapsto y_0$. Let
$q$ be the class of $ax_0^2 + bx_0y_0 + cy_0^2$ in
$\text{Sym}^2_\kappa(\mathfrak m/\mathfrak m^2)$.
Write $Q(x, y) = ax^2 + bxy + cy^2$ viewed as a polynomial
in two variables. Then we see that
$$
Q(x_0, y_0) = ax_0^2 + bx_0y_0 + cy_0^2 =
\sum\nolimits_{i + j = 3} a_{ij} x_0^iy_0^j
$$
for some $a_{ij}$ in $A$. We want to prove that we can
increase the order of vanishing by changing our choice
of $x_0$, $y_0$. Suppose that $x_1, y_1 \in \mathfrak m^2$.
Then
$$
Q(x_0 + x_1, y_0 + y_1) = Q(x_0, y_0) +
(2ax_0 + by_0)x_1 + (bx_0 + 2cy_0)y_1 \bmod \mathfrak m^4
$$
Nondegeneracy of $Q$ means exactly that $2ax_0 + by_0$ and $bx_0 + 2cy_0$
are a $\kappa$-basis for $\mathfrak m/\mathfrak m^2$, see discussion
preceding the lemma. Hence we can
certainly choose $x_1, y_1 \in \mathfrak m^2$ such that
$Q(x_0 + x_1, y_0 + y_1) \in \mathfrak m^4$.
Continuing in this fashion by induction
we can find $x_i, y_i \in \mathfrak m^{i + 1}$
such that
$$
Q(x_0 + x_1 + \ldots + x_n, y_0 + y_1 + \ldots + y_n) \in \mathfrak m^{n + 3}
$$
Since $A$ is complete we can set
$x_\infty = \sum x_i$ and $y_\infty = \sum y_i$
and we can consider the map $\kappa[[x, y]] \longrightarrow A$
sending $x$ to $x_\infty$ and $y$ to $y_\infty$. This map
induces a surjection $\kappa[[x, y]]/(Q) \longrightarrow A$ by
Algebra, Lemma \ref{algebra-lemma-completion-generalities}.
By Lemma \ref{lemma-reduced-quotient-regular-ring-dim-2}
the kernel of $k[[x, y]] \to A$ is principal.
But the kernel cannot contain a proper divisor of $Q$
as such a divisor would have degree $1$ in $x, y$
and this would contradict $\dim(\mathfrak m/\mathfrak m^2) = 2$.
Hence $Q$ generates the kernel as desired.
\end{proof}

\end{document}
